\documentclass[11pt]{article}

\usepackage[a4paper,margin=19mm,headheight=24pt]{geometry}
\usepackage{fontspec}
\usepackage{xcolor}
\usepackage{booktabs}
\usepackage{array}
\usepackage{enumitem}
\usepackage{fancyhdr}
\usepackage{hyperref}
\usepackage{listings}
\usepackage{amsmath}

\setmainfont{DejaVu Sans}
\setsansfont{DejaVu Sans}
\setmonofont{DejaVu Sans Mono}[Scale=0.86]

\definecolor{navy}{HTML}{16324F}
\definecolor{blue}{HTML}{2474A6}
\definecolor{cyan}{HTML}{EAF5FA}
\definecolor{green}{HTML}{237A57}
\definecolor{palegreen}{HTML}{EAF6F0}
\definecolor{amber}{HTML}{A86F16}
\definecolor{paleamber}{HTML}{FFF6E5}
\definecolor{ink}{HTML}{24313B}
\definecolor{muted}{HTML}{62717D}
\definecolor{rule}{HTML}{D5DEE5}

\hypersetup{
  colorlinks=true,
  linkcolor=blue,
  urlcolor=blue,
  pdfauthor={RICH/STORM engineering note},
  pdftitle={Reducing Monte Carlo Buffer Memory Without Sacrificing Performance}
}

\pagestyle{fancy}
\fancyhf{}
\lhead{\textcolor{muted}{RICH / STORM Engineering Note}}
\rhead{\textcolor{muted}{Monte Carlo Buffer Memory}}
\cfoot{\textcolor{muted}{\thepage}}
\renewcommand{\headrulewidth}{0.4pt}
\renewcommand{\headrule}{\hbox to\headwidth{\color{rule}\leaders\hrule height \headrulewidth\hfill}}

\setlength{\parindent}{0pt}
\setlength{\parskip}{6pt}
\setlength{\emergencystretch}{2em}
\setlist[itemize]{leftmargin=5mm,itemsep=3pt,topsep=3pt}
\setlist[enumerate]{leftmargin=7mm,itemsep=5pt,topsep=4pt}

\lstset{
  basicstyle=\ttfamily\small,
  backgroundcolor=\color{cyan},
  frame=single,
  rulecolor=\color{rule},
  framesep=7pt,
  xleftmargin=2pt,
  xrightmargin=2pt,
  columns=fullflexible,
  keepspaces=true,
  showstringspaces=false
}

\newcommand{\sectiontitle}[1]{%
  \vspace{4pt}%
  {\Large\bfseries\color{navy} #1}\par
  \vspace{2pt}\textcolor{rule}{\rule{\linewidth}{0.8pt}}\vspace{3pt}
}

\newcommand{\callout}[3]{%
  \noindent\colorbox{#1}{%
    \parbox{\dimexpr\linewidth-2\fboxsep}{%
      \textcolor{#2}{\bfseries #3}%
    }%
  }%
}

\begin{document}

\begin{center}
  \vspace*{8mm}
  {\fontsize{25}{29}\selectfont\bfseries\color{navy}
    Reducing Monte Carlo Buffer Memory\par}
  \vspace{2mm}
  {\fontsize{18}{22}\selectfont\color{blue}
    without sacrificing transport performance\par}
  \vspace{7mm}
  {\large RICH / STORM --- \texttt{moving\_slab\_mc\_32}}\par
  \vspace{2mm}
  {\color{muted}Source-level engineering assessment \textbullet{} 9 August 2026}
\end{center}

\vspace{8mm}
\callout{cyan}{navy}{Executive finding}

The particle representation should remain unchanged. The largest safe memory
reduction comes from eliminating unnecessary capacity and redundant full-buffer
copies. A bulk append currently allocates 50\% more handler capacity than the
known required size while the complete generation vector is still alive.

\vspace{2mm}
\begin{center}
\renewcommand{\arraystretch}{1.45}
\begin{tabular}{>{\raggedright\arraybackslash}p{0.25\linewidth}
                >{\centering\arraybackslash}p{0.18\linewidth}
                >{\centering\arraybackslash}p{0.20\linewidth}
                >{\raggedright\arraybackslash}p{0.25\linewidth}}
\toprule
\textbf{Stage} & \textbf{Generation vector} & \textbf{Handler capacity} & \textbf{Particle-slot peak} \\
\midrule
Current & $N$ & $1.5N$ & $2.5N$ \\
Corrected growth & $N$ & $N$ for a large bulk append & $2N$ \\
Direct generation & none & approximately $N$ & approximately $N$ \\
\bottomrule
\end{tabular}
\end{center}

For the current 208-byte particle, the first-generation lower-bound peak is
approximately $2.5N\times208=520N$ bytes before allocator, MPI, mesh, and physics
overheads. The problem is therefore not explained by particle payload alone.

\sectiontitle{1. Primary fix: correct the capacity equation}

The bulk-append logic in
\path{source/monte/manager/parallel/RankHandler2.hpp:384} effectively computes:

\begin{lstlisting}
new_capacity = max(1.5 * old_capacity, 1.5 * required);
\end{lstlisting}

Use standard amortized growth instead:

\begin{lstlisting}
new_capacity = max(1.5 * old_capacity, required);
\end{lstlisting}

This preserves 1.5x growth for incremental arrivals, so repeated small appends
retain amortized constant-time behavior. A large append whose final size is
already known receives exactly the capacity it requires rather than an unused
50\% surplus. The equivalent remote-growth calculation near
\path{RankHandler2.hpp:704} should follow the same rule.

\callout{palegreen}{green}{Why this should not hurt performance}

The number of reallocations for incremental traffic remains governed by the
same 1.5x policy. Large bulk operations allocate less memory, initialize fewer
pages, and register a smaller RDMA region. This should reduce both memory use
and allocation time.

\sectiontitle{2. Remove duplicate copies during resize}

The current resize path first copies every live particle into
\texttt{activeParticles} at \path{RankHandler2.hpp:455}. It then allocates the
new particle array and copies the temporary data again.

For the self handler, copy the circular queue directly into the new array in
logical order, then replace the old allocation. This removes:

\begin{itemize}
  \item one live-particle-sized temporary allocation;
  \item one complete memory copy;
  \item the corresponding peak RSS and memory-bandwidth cost.
\end{itemize}

The native IBV path contains a second redundancy. The relevant implementation is
{\small\path{source/utils/rma/IBVRemoteMemoryAgent.hpp:520}.}
\texttt{LocalResize} copies the entire old capacity into an internal
\texttt{saved} vector even
though \texttt{RankHandler2} has already preserved the live particles. A
discarding or compacting resize operation should preserve the queue exactly
once, release the old registration, allocate/register the new region, and
restore only the live slots.

\sectiontitle{3. Reclaim cold high-water capacities with hysteresis}

Do not shrink buffers during the transport hot loop. Reclaim memory only at a
step boundary, when the manager already verifies that handler queues are empty.

Recommended policy:

\begin{enumerate}
  \item Track each buffer's recent peak occupancy over several steps.
  \item Shrink only after multiple cold steps and only when capacity exceeds
        roughly twice the recent peak.
  \item Retain modest headroom, for example 1.25x the recent peak.
  \item Coordinate peer-handler resizing only at an existing synchronization
        point.
\end{enumerate}

The current shrink logic in
\path{source/monte/manager/parallel/RDMAMonteCarloManager.hpp:1532} excludes
the self handler, normally runs only every 50 steps, and shrinks neighbor
handlers only under restrictive conditions. A historical spike can therefore
remain resident for the rest of the run.

Send buffers and detached-particle vectors also use \texttt{clear()}, which
retains their peak capacity. Apply the same cold-buffer policy only to unusually
large retained allocations. Avoid \texttt{shrink\_to\_fit()} after every flush;
that would create allocation churn and harm performance.

\sectiontitle{4. Larger follow-up: generate directly into final storage}

The generated-particle vector remains alive while its particles are copied into
the self handler. The relevant sequence begins near
\path{RDMAMonteCarloManager.hpp:2174}.

A higher-value follow-up is to count the particles first, reserve the exact
self-handler span, and let each cell or worker generate into a disjoint range.
This removes the full \texttt{newParticles1} allocation and the subsequent copy.
It should improve performance through fewer allocations and less memory traffic,
but it requires a broader physics/manager API change than the growth correction.

\sectiontitle{5. Recommended implementation order}

\renewcommand{\arraystretch}{1.45}
\small
\begin{tabular}{>{\centering\arraybackslash}p{0.07\linewidth}
                >{\raggedright\arraybackslash}p{0.50\linewidth}
                >{\raggedright\arraybackslash}p{0.31\linewidth}}
\toprule
\textbf{Order} & \textbf{Change} & \textbf{Expected effect} \\
\midrule
1 & Correct bulk growth to $\max(1.5\,\text{old},\text{required})$ & Large immediate reduction; lowest risk \\
2 & Copy the circular queue directly during resize & Lower transient peak; faster resize \\
3 & Add IBV/OFI compact-or-discard resize & Remove redundant full-capacity save/copy \\
4 & Add cold-buffer hysteresis for self, peer, send, and detached buffers & Release historical spikes without hot-loop churn \\
5 & Generate directly into pre-sized final storage & Largest long-term peak reduction \\
\bottomrule
\end{tabular}
\normalsize

\vspace{4mm}
\callout{paleamber}{amber}{Performance guardrails}

Keep RDMA transfer batches large. Keep geometric growth for incremental traffic.
Do not shrink active buffers or shrink on every flush. Validate using capacity,
live-occupancy, resize-count, bytes-copied, transport-time, and peak-RSS
diagnostics per rank and per node.

\vfill
{\footnotesize\color{muted}
This note is based on a static audit of the current source tree. No RICH build or
runtime validation was performed while preparing it.}

\end{document}
